El teorema anterior describe qué ocurre al cambiar la masa. También necesitamos
comparar dos posiciones de un mismo tablero. Una cota Lipschitz controla el cambio
del perfil por la distancia entre esas posiciones. La cuestión es si la misma
constante sigue sirviendo cuando quedan cada vez menos filas vacías.

\begin{teorema}[Regularidad espacial]\label{thm:lipschitz_intro}
\begin{enumerate}[(a)]
\item Para $\rho\in(0,1)$ y $\lambda=\lambda(\rho)$, las vacancias y las proyecciones son Lipschitz en
sus intervalos cerrados con constante $L(\lambda)=2\lambda+8\lambda^2(1+\lambda)$, y $\nu_\rho$ es
Lipschitz en el cuadrado con constante $6\lambda L(\lambda)$ para la norma del máximo. Las cotas son
uniformes en cada compacto de $(0,1)$.
\item Existe $L_*<\infty$ tal que, para todo $x\in[0,1/2]$, los factores reescalados
$r_R,r_C,u_D,u_A$ son Lipschitz con constante $L_*$, y $\nu$ es acotada y Lipschitz con constantes
independientes de $x$.
\end{enumerate}
\end{teorema}

La parte (a) degenera cuando $\rho\to1$, porque $\lambda(\rho)\to\infty$; la parte (b) es la que
permite trabajar hasta el tablero lleno.

Definimos, para $x\in(0,1)$ y $\rho=1-x$,
\begin{equation}\label{eq:F}
 F(x)=\frac{1-x}{\lambda(1-x)\,x^2}=\frac{1-x}{b}=\int_{[0,1]^2}r_R(r)\,r_C(c)\,u_D(r-c)\,u_A(r+c)\,dr\,dc .
\end{equation}
La última expresión es la integral de $u_D(r-c)\,u_A(r+c)$ respecto del peso $r_R(r)\,r_C(c)$,
que es una densidad de probabilidad en el cuadrado. El interés de $F$ está en la identidad
integral del teorema siguiente, que recupera $\sopt(1)$ a partir de ella.

\begin{teorema}[La función $F$ y el valor óptimo]\label{thm:F_intro}
La función $F$ se extiende a una función real analítica y positiva en un entorno de $[0,1]$, con
\[
 F(0)=\frac1{b_0},\qquad F(1)=1,\qquad \int_0^1\log F(x)\,dx=1+\sopt(1).
\]
Además, cuando $x\downarrow0$,
\[
 \sopt(1-x)-\sopt(1)=2x\log\frac1x+O(x).
\]
\end{teorema}

\begin{observacion}[Valores numéricos]
Una resolución numérica determinista del sistema de Euler (Sección~\ref{sec:discusion}) da
$F(0)\approx0{,}0809$, $F'(0)\approx0{,}516$ e $\int_0^1\log F\approx-0{,}9440$, de modo que
$\sopt(1)\approx-1{,}9440$. Este valor es compatible con los valores conocidos de la constante
$\alpha$~\cite{Simkin,NAB}, como cabe esperar si $\alpha=-\sopt(1)$
(véase la Observación~\ref{obs:simkin}). Son valores numéricos, sin
certificación de error.
\end{observacion}

\subsection{Del maximizador al perfil límite}

Para la masa uno, la existencia y la unicidad del óptimo de la entropía
figuran en el Lema~2.22 de Simkin~\cite{Simkin}. La
Observación~\ref{obs:simkin} precisa la correspondencia de los funcionales.
La familia de masas parciales permite estudiar cómo se llega a ese óptimo.
Primero deducimos la ecuación de Euler en el interior (Sección~\ref{sec:euler}).
Sus cotas reescaladas permiten identificar los factores que subsisten en masa uno
(Secciones~\ref{sec:cotas} y~\ref{sec:limite}). El análisis de la dependencia
respecto de la masa se completa en la Sección~\ref{sec:analiticidad}.
Para controlar los cambios de posición comparamos el perfil con sus traslaciones
(Sección~\ref{sec:regularidad}). Finalmente, al integrar la ecuación de Euler
se obtiene la relación entre $F$ y la entropía óptima (Sección~\ref{sec:F}),
que evaluamos numéricamente en la Sección~\ref{sec:discusion}.

\subsection{Relación con trabajos previos}

Simkin~\cite{Simkin} acota el valor óptimo en masa uno discretizando el problema; su cota
superior usa la dualidad de Lagrange, que da al maximizador discreto una forma exponencial en los
multiplicadores~\cite[Sección~6]{Simkin}. Nobel, Agrawal y Boyd~\cite{NAB} resuelven una
discretización del mismo problema con un método de Newton y obtienen el intervalo citado. Que el
maximizador de una entropía bajo restricciones lineales tenga forma exponencial es un principio
clásico~\cite{Csiszar}. Lo que añade este trabajo es el estudio de la familia de masas parciales,
con su ecuación de Euler, y sobre todo el control uniforme hasta el tablero lleno: cotas
explícitas, convergencia y analiticidad del perfil reescalado, regularidad Lipschitz uniforme y
la identidad integral para $F$. Hasta donde sabemos, estos resultados son nuevos; no se reclama novedad para las técnicas generales de
análisis convexo, de dualidad y de la función implícita~\cite{Brezis,Zeidler}.
